\documentclass[UTF8,12pt]{ctexart}
\usepackage[a4paper,margin=24mm]{geometry}
\usepackage{amsmath,amssymb,bm,graphicx,adjustbox}
\newcommand{\fitmath}[1]{\begin{adjustbox}{max width=\linewidth}$\displaystyle #1$\end{adjustbox}}
\setlength{\parindent}{0pt}
\setlength{\parskip}{.7em}
\sloppy
\allowdisplaybreaks
\begin{document}
\section*{2013 年考研数学三 · 真题}
\subsection*{原 PDF 第 31 页}

\subsection*{2013 年全国硕士研究生招生考试试题}

\subsection*{一、选择题（本题共 8 小题，每小题 4 分，共 32 分。在每小题给出的四个选项中，只有一项符合题目要求，把所选项前的字母填在题后的括号内。）}

(1) 当 \(\displaystyle x\to0\) 时，用“\(\displaystyle o(x)\)”表示比 \(\displaystyle x\) 高阶的无穷小量，则下列式子中错误的是（　）。


(A) \(\displaystyle x\cdot o(x^2)=o(x^3)\)。 (B) \(\displaystyle o(x)\cdot o(x^2)=o(x^3)\)。 (C) \(\displaystyle o(x^2)+o(x^2)=o(x^2)\)。 (D) \(\displaystyle o(x)+o(x^2)=o(x^2)\)。


(2) 函数 \(\displaystyle f(x)=\dfrac{\displaystyle |x|^x-1}{\displaystyle x(x+1)\ln|x|}\) 的可去间断点的个数为（　）。


(A) 0。 (B) 1。 (C) 2。 (D) 3。


(3) 设 \(\displaystyle D_k\) 是圆域 \(\displaystyle D=\{(x,\allowbreak y)\mid x^2+y^2\leq1\}\) 位于第 \(\displaystyle k\) 象限的部分。记 \(\displaystyle I_k=\displaystyle \iint_{D_k}(y-x)\,dx\,dy\)（\(\displaystyle k=1,\allowbreak 2,\allowbreak 3,\allowbreak 4\)），则（　）。


(A) \(\displaystyle I_1>0\)。 (B) \(\displaystyle I_2>0\)。 (C) \(\displaystyle I_3>0\)。 (D) \(\displaystyle I_4>0\)。


(4) 设 \(\displaystyle \{a_n\}\) 为正项数列，下列选项正确的是（　）。


(A) 若 \(\displaystyle a_n>a_{n+1}\)，则 \(\displaystyle \sum_{n=1}^{\infty}(-1)^{n-1}a_n\) 收敛。


(B) 若 \(\displaystyle \sum_{n=1}^{\infty}(-1)^{n-1}a_n\) 收敛，则 \(\displaystyle a_n>a_{n+1}\)。


(C) 若 \(\displaystyle \sum_{n=1}^{\infty}a_n\) 收敛，则存在常数 \(\displaystyle p>1\)，使 \(\displaystyle \lim_{n\to\infty}n^pa_n\) 存在。


(D) 若存在常数 \(\displaystyle p>1\)，使 \(\displaystyle \lim_{n\to\infty}n^pa_n\) 存在，则 \(\displaystyle \sum_{n=1}^{\infty}a_n\) 收敛。


(5) 设 \(\displaystyle A,\allowbreak B,\allowbreak C\) 均为 \(\displaystyle n\) 阶矩阵。若 \(\displaystyle AB=C\)，且 \(\displaystyle B\) 可逆，则（　）。


(A) 矩阵 C 的行向量组与矩阵 A 的行向量组等价。 (B) 矩阵 C 的列向量组与矩阵 A 的列向量组等价。 (C) 矩阵 C 的行向量组与矩阵 B 的行向量组等价。 (D) 矩阵 C 的列向量组与矩阵 B 的列向量组等价。


(6) 矩阵 \(\displaystyle \begin{pmatrix}1&a&1\\a&b&a\\1&a&1\end{pmatrix}\) 与 \(\displaystyle \begin{pmatrix}2&0&0\\0&b&0\\0&0&0\end{pmatrix}\) 相似的充分必要条件为（　）。


(A) \(\displaystyle a=0,\allowbreak b=2\)。 (B) \(\displaystyle a=0,\allowbreak b\) 为任意常数。 (C) \(\displaystyle a=2,\allowbreak b=0\)。 (D) \(\displaystyle a=2,\allowbreak b\) 为任意常数。


(7) 设 \(\displaystyle X_1,\allowbreak X_2,\allowbreak X_3\) 是随机变量，且 \(\displaystyle X_1\sim N(0,\allowbreak 1),\allowbreak \ X_2\sim N(0,\allowbreak 2^2),\allowbreak \ X_3\sim N(5,\allowbreak 3^2)\)，\(\displaystyle p_i=P\{-2\leq X_i\leq2\}\)（\(\displaystyle i=1,\allowbreak 2,\allowbreak 3\)），则（　）。


(A) \(\displaystyle p_1>p_2>p_3\)。 (B) \(\displaystyle p_2>p_1>p_3\)。 (C) \(\displaystyle p_3>p_1>p_2\)。 (D) \(\displaystyle p_1>p_3>p_2\)。


\clearpage

\subsection*{原 PDF 第 32 页}

(8) 设随机变量 \(\displaystyle X\) 和 \(\displaystyle Y\) 相互独立，且 \(\displaystyle X\) 和 \(\displaystyle Y\) 的概率分布分别为：\(\displaystyle \begin{array}{c|cccc}X&0&1&2&3\\\hline P&\dfrac{\displaystyle 1}{\displaystyle 2}&\dfrac{\displaystyle 1}{\displaystyle 4}&\dfrac{\displaystyle 1}{\displaystyle 8}&\dfrac{\displaystyle 1}{\displaystyle 8}\end{array}\)，\(\displaystyle \begin{array}{c|ccc}Y&-1&0&1\\\hline P&\dfrac{\displaystyle 1}{\displaystyle 3}&\dfrac{\displaystyle 1}{\displaystyle 3}&\dfrac{\displaystyle 1}{\displaystyle 3}\end{array}\)。则 \(\displaystyle P\{X+Y=2\}=\)（　）。


(A) \(\displaystyle \dfrac{\displaystyle 1}{\displaystyle 12}\)。 (B) \(\displaystyle \dfrac{\displaystyle 1}{\displaystyle 8}\)。 (C) \(\displaystyle \dfrac{\displaystyle 1}{\displaystyle 6}\)。 (D) \(\displaystyle \dfrac{\displaystyle 1}{\displaystyle 2}\)。


\subsection*{二、填空题（本题共 6 小题，每小题 4 分，共 24 分，把答案填在题中横线上。）}

(9) 设曲线 \(\displaystyle y=f(x)\) 与 \(\displaystyle y=x^2-x\) 在点 \(\displaystyle (1,\allowbreak 0)\) 处有公共切线，则 \(\displaystyle \lim_{n\to\infty}nf\left(\dfrac{\displaystyle n}{\displaystyle n+2}\right)=\)\_\_\_\_\_\_。


(10) 设函数 \(\displaystyle z=z(x,\allowbreak y)\) 由方程 \(\displaystyle (z+y)^x=xy\) 确定，则 \(\displaystyle \left.\dfrac{\displaystyle \partial z}{\displaystyle \partial x}\right|_{(1,\allowbreak 2)}=\)\_\_\_\_\_\_。


(11) \(\displaystyle \int_1^{+\infty}\dfrac{\displaystyle \ln x}{\displaystyle (1+x)^2}\,dx=\)\_\_\_\_\_\_。


(12) 微分方程 \(\displaystyle y^{\prime\prime}-y^{\prime}+\dfrac{\displaystyle 1}{\displaystyle 4}y=0\) 的通解为 \(\displaystyle y=\)\_\_\_\_\_\_。


(13) 设 \(\displaystyle A=(a_{ij})\) 是 3 阶非零矩阵，\(\displaystyle |A|\) 为 \(\displaystyle A\) 的行列式，\(\displaystyle A_{ij}\) 为 \(\displaystyle a_{ij}\) 的代数余子式。若 \(\displaystyle a_{ij}+A_{ij}=0\)（\(\displaystyle i,\allowbreak j=1,\allowbreak 2,\allowbreak 3\)），则 \(\displaystyle |A|=\)\_\_\_\_\_\_。


(14) 设随机变量 \(\displaystyle X\) 服从标准正态分布 \(\displaystyle N(0,\allowbreak 1)\)，则 \(\displaystyle E(Xe^{2X})=\)\_\_\_\_\_\_。


\subsection*{三、解答题（本题共 9 小题，共 94 分，解答应写出文字说明、证明过程或演算步骤。）}

(15)（本题满分 10 分）当 \(\displaystyle x\to0\) 时，\(\displaystyle 1-\cos x\cdot\cos2x\cdot\cos3x\) 与 \(\displaystyle ax^n\) 为等价无穷小量，求 \(\displaystyle n\) 与 \(\displaystyle a\) 的值。


(16)（本题满分 10 分）设 \(\displaystyle D\) 是由曲线 \(\displaystyle y=x^{1/3}\)，直线 \(\displaystyle x=a\)（\(\displaystyle a>0\)）及 \(\displaystyle x\) 轴所围成的平面图形，\(\displaystyle V_x,\allowbreak V_y\) 分别是 \(\displaystyle D\) 绕 \(\displaystyle x\) 轴、\(\displaystyle y\) 轴旋转一周所得旋转体的体积。若 \(\displaystyle V_y=10V_x\)，求 \(\displaystyle a\) 的值。


\clearpage

\subsection*{原 PDF 第 33 页}

(17)（本题满分 10 分）设平面区域 \(\displaystyle D\) 由直线 \(\displaystyle x=3y,\allowbreak \ y=3x\) 及 \(\displaystyle x+y=8\) 围成，计算 \(\displaystyle \iint_D x^2\,dx\,dy\)。


(18)（本题满分 10 分）设生产某商品的固定成本为 60 000 元，可变成本为 20 元／件，价格函数为 \(\displaystyle p=60-\dfrac{\displaystyle Q}{\displaystyle 1000}\)（\(\displaystyle p\) 是单价，单位：元；\(\displaystyle Q\) 是销量，单位：件）。已知产销平衡，求：


（I）该商品的边际利润；（II）当 \(\displaystyle p=50\) 时的边际利润，并解释其经济意义；（III）使得利润最大的定价 \(\displaystyle p\)。


(19)（本题满分 10 分）设函数 \(\displaystyle f(x)\) 在 \(\displaystyle [0,\allowbreak +\infty)\) 上可导，\(\displaystyle f(0)=0\) 且 \(\displaystyle \lim_{x\to+\infty}f(x)=2\)。证明：


（I）存在 \(\displaystyle a>0\)，使得 \(\displaystyle f(a)=1\)；（II）对（I）中的 \(\displaystyle a\)，存在 \(\displaystyle \xi\in(0,\allowbreak a)\)，使得 \(\displaystyle f^{\prime}(\xi)=\dfrac{\displaystyle 1}{\displaystyle a}\)。


(20)（本题满分 11 分）设 \(\displaystyle A=\begin{pmatrix}1&a\\1&0\end{pmatrix},\allowbreak \ B=\begin{pmatrix}0&1\\1&b\end{pmatrix}\)。当 \(\displaystyle a,\allowbreak b\) 为何值时，存在矩阵 \(\displaystyle C\) 使得 \(\displaystyle AC-CA=B\)，并求所有矩阵 \(\displaystyle C\)。


\clearpage

\subsection*{原 PDF 第 34 页}

(21)（本题满分 11 分）设二次型 \(\displaystyle f(x_1,\allowbreak x_2,\allowbreak x_3)=2(a_1x_1+a_2x_2+a_3x_3)^2+(b_1x_1+b_2x_2+b_3x_3)^2\)，记 \(\displaystyle \alpha=\begin{pmatrix}a_1\\a_2\\a_3\end{pmatrix},\allowbreak \ \beta=\begin{pmatrix}b_1\\b_2\\b_3\end{pmatrix}\)。


（I）证明二次型 \(\displaystyle f\) 对应的矩阵为 \(\displaystyle 2\alpha\alpha^{\mathrm T}+\beta\beta^{\mathrm T}\)；（II）若 \(\displaystyle \alpha,\allowbreak \beta\) 正交且均为单位向量，证明 \(\displaystyle f\) 在正交变换下的标准形为 \(\displaystyle 2y_1^2+y_2^2\)。


(22)（本题满分 11 分）设 \(\displaystyle (X,\allowbreak Y)\) 是二维随机变量，\(\displaystyle X\) 的边缘概率密度为 \(\displaystyle f_X(x)=\begin{cases}3x^2,\allowbreak &0<x<1,\allowbreak \\0,\allowbreak &\text{其他}\end{cases}\)。在给定 \(\displaystyle X=x\)（\(\displaystyle 0<x<1\)）的条件下 \(\displaystyle Y\) 的条件概率密度为 \(\displaystyle f_{Y|X}(y|x)=\begin{cases}\dfrac{\displaystyle 3y^2}{\displaystyle x^3},\allowbreak &0<y<x,\allowbreak \\0,\allowbreak &\text{其他}\end{cases}\)。


（I）求 \(\displaystyle (X,\allowbreak Y)\) 的概率密度 \(\displaystyle f(x,\allowbreak y)\)；（II）求 \(\displaystyle Y\) 的边缘概率密度 \(\displaystyle f_Y(y)\)；（III）求 \(\displaystyle P\{X>2Y\}\)。


(23)（本题满分 11 分）设总体 \(\displaystyle X\) 的概率密度为 \(\displaystyle f(x;\theta)=\begin{cases}\dfrac{\displaystyle \theta^2}{\displaystyle x^3}e^{-\theta/x},\allowbreak &x>0,\allowbreak \\0,\allowbreak &\text{其他},\allowbreak \end{cases}\) 其中 \(\displaystyle \theta\) 为未知参数且大于零。\(\displaystyle X_1,\allowbreak X_2,\allowbreak \ldots,\allowbreak X_n\) 为来自总体 \(\displaystyle X\) 的简单随机样本。


（I）求 \(\displaystyle \theta\) 的矩估计量；（II）求 \(\displaystyle \theta\) 的最大似然估计量。


\clearpage
\end{document}
